% Πίνακας περιεχομένων: ίδια ιεραρχία σε PRINT και SCREEN, μόνο μαύρο.
% Οι άτιτλες ως προς την αρίθμηση ομάδες/εγγραφές έχουν δικούς τους τύπους:
% η επόμενη γραμμή τους ευθυγραμμίζεται με την πρώτη, χωρίς κενή θέση αριθμού.
\makeatletter
\newcommand*{\toclevel@bookcollection}{1}
\newcommand*{\toclevel@bookitem}{2}
\newcommand*{\l@bookcollection}[2]{%
  \ifnum\c@tocdepth>0\relax
    % Μένει μαζί με τουλάχιστον την αρχή της πρώτης καταχώρισης.
    \Needspace{3\baselineskip}%
    \@dottedtocline{1}{1.5em}{0em}{#1}{#2}%
    \nobreak
  \fi
}
\newcommand*{\l@bookitem}{\@dottedtocline{2}{3.8em}{0em}}
% Χώρος για τριψήφιους αριθμούς και σαφές διάστημα από μακρύ τίτλο.
\renewcommand{\@pnumwidth}{2em}
\renewcommand{\@tocrmarg}{3em}
% Ο σελιδοδείκτης δημιουργείται ΜΕΣΑ στον ήδη υπάρχοντα τίτλο:
% δείχνει στην κορυφή του τίτλου και δεν προκαλεί νέα/κενή σελίδα.
\patchcmd{\tableofcontents}
  {\contentsname}
  {\pdfbookmark[0]{\contentsname}{book.contents}\contentsname}
  {}{\PackageError{universalbook}{Cannot attach Contents bookmark}{The book-class tableofcontents definition changed.}}
\makeatother
% Ενεργός σύνδεσμος τόσο στον τίτλο όσο και στον αριθμό σελίδας.
\hypersetup{linktoc=all}
